\begin{enumerate}
\item آمین در موقعیت ۷ هسته سفالوسپورین به کربونیل فعال مشتق ماندلیک حمله می‌کند و حدواسط \lr{tetrahedral intermediate} تشکیل می‌شود.
\item حدواسط فرو می‌ریزد، حلقه کربنات باز می‌شود و با دکربوکسیلاسیون $\mathrm{CO_2}$ آزاد می‌گردد.
\item انتقال پروتون، آمید پایدار یعنی سفاماندول را می‌دهد.
\end{enumerate}

\schemestart
\chemfig{R-NH_2}
\+
\chemfig{Ph-CH(-OH)-C(=O)-O-[:0]*5(O-C(=O)-O-C(=O)-O-)}
\arrow{->[\scriptsize حمله نوکلئوفیلی]}
\chemfig{R-NH-C(OH)(O-[:0]*5(O-C(=O)-O-C(=O)-O-))-CH(-OH)-Ph}
\arrow{->[\scriptsize فروپاشی و خروج $\mathrm{CO_2}$]}
\chemfig{R-NH-C(=O)-CH(-OH)-Ph}
\+
\chemfig{CO_2}
\schemestop

\smallskip
در این شماتیک، $R$ نماینده هسته سفالوسپورینی با گروه تترازول در موقعیت ۳ و گروه کربوکسیلیک در موقعیت ۴ است.